% from ega3-1-fr.tex
\oldpage[III]{137}

\begin{corollary}[4.4.6]
\label{III.4.4.6-fr}
Soient $A$ un anneau noethérien, $B$ une $A$-algèbre de type fini,
$\mathfrak q$ un idéal premier de $B$, $\mathfrak p$ son image réciproque
dans $A$. On suppose que $\mathfrak q$ soit à la fois maximal et minimal dans
l'ensemble des idéaux premiers de $B$ dont l'image réciproque est
$\mathfrak p$. Alors il existe $g\in B-\mathfrak q$, une $A$-algèbre finie
$A'$ et un élément $f'\in A'$ tels que les $A$-algèbres $B_g$ et $A'_{f'}$
soient isomorphes.
\end{corollary}

\begin{proof}
Il suffit en effet d'appliquer (\hyperref[III.4.4.5-fr]{4.4.5}) à
$Y=\operatorname{Spec}(A)$ et $X=\operatorname{Spec}(B)$, l'hypothèse sur
$\mathfrak q$ signifiant exactement qu'il est isolé dans sa fibre
(I, 1.1.7).
\end{proof}

On en déduit le résultat suivant, moins général en apparence~:

\begin{corollary}[4.4.7]
\label{III.4.4.7-fr}
Soient $A$ un anneau local noethérien, $B$ une $A$-algèbre de type fini,
$\mathfrak n$ un idéal premier de $B$ dont l'image réciproque dans $A$ est
l'idéal maximal $\mathfrak m$. On suppose que $\mathfrak n$ est maximal dans
$B$ et est minimal dans l'ensemble des idéaux premiers de $B$ dont l'image
réciproque est $\mathfrak m$ (ce qui signifie aussi que
$B\mathfrak m$ est primaire pour $\mathfrak n$). Alors il existe une
$A$-algèbre finie $A'$ et un idéal maximal $\mathfrak m'$ de $A'$ (dont
$\mathfrak m$ est l'image réciproque dans $A$) tels que $B_{\mathfrak n}$
soit isomorphe à la $A$-algèbre $A'_{\mathfrak m'}$.
\end{corollary}

Le cas particulier suivant de (\hyperref[III.4.4.7-fr]{4.4.7}) est aussi
parfois appelé «~Main Theorem~»~:

\begin{corollary}[4.4.8]
\label{III.4.4.8-fr}
Sous les conditions de (\hyperref[III.4.4.7-fr]{4.4.7}), supposons en outre
$A$ et $B$ intègres et ayant même corps des fractions $K$. Alors, si $A$ est
intégralement clos, on a $B=A$.
\end{corollary}

\begin{proof}
En effet, la Remarque (\hyperref[III.4.4.4-fr]{4.4.4}) montre que l'on peut
supposer, dans l'application de (\hyperref[III.4.4.7-fr]{4.4.7}) que $A'$ est
intègre et a $K$ pour corps des fractions~; l'hypothèse sur $A$ entraîne alors
$A'=A$, donc $B_{\mathfrak n}=A$~; comme on a
$A\subset B\subset B_{\mathfrak n}$, on en conclut bien $A=B$.
\end{proof}

L'énoncé (\hyperref[III.4.4.8-fr]{4.4.8}) est la forme donnée par Zariski à
son «~Main theorem~» (étendu aux anneaux locaux noethériens intègres
quelconques).

Les corollaires précédents étaient des variantes de nature locale de
(\hyperref[III.4.4.3-fr]{4.4.3}), qui est un résultat global. Voici une autre
conséquence de nature globale~:

\begin{corollary}[4.4.9]
\label{III.4.4.9-fr}
Soient $Y$ un préschéma intègre localement noethérien, $f:X\to Y$ un
morphisme séparé, de type fini et birationnel. Supposons en outre $Y$ normal
et toutes les fibres $f^{-1}(y)$ finies pour $y\in Y$. Alors $f$ est une
immersion ouverte~; si en outre $f$ est fermé (et en particulier si $f$ est
propre), $f$ est un isomorphisme.
\end{corollary}

\begin{proof}
Soit en effet $x\in X$, et posons $y=f(x)$. Comme $f^{-1}(y)$ est un schéma
algébrique sur $k(y)$, l'hypothèse qu'il est fini entraîne qu'il est discret
(I, 6.4.4)~; en outre $\mathcal O_y$ est intégralement clos et
$\mathcal O_x$ et $\mathcal O_y$ ont même corps des fractions (I, 7.1.5).
On peut donc appliquer (\hyperref[III.4.4.8-fr]{4.4.8}), et si
$f=(\psi,\theta)$, l'homomorphisme
$\theta_y^\#:\mathcal O_y\to\mathcal O_x$ est bijectif~; on en conclut
(I, 6.5.4) que $f$ est un isomorphisme local. Mais comme $f$ est séparé et
$X$ intègre, $f$ est une immersion ouverte (I, 8.2.8). La dernière assertion
résulte de ce que $f$ est dominant.
\end{proof}

\begin{proposition}[4.4.10]
\label{III.4.4.10-fr}
Soient $Y$ un préschéma localement noethérien, $f:X\to Y$ un morphisme
localement de type fini. L'ensemble $X'$ des $x\in X$ isolés dans leur fibre
$f^{-1}(f(x))$ est ouvert dans $X$.
\end{proposition}

\begin{proof}
La question étant locale sur $X$ et $Y$, on peut supposer $X$ et $Y$ affines
noethériens et $f$ de type fini~; $f$ est alors un morphisme affine de type
fini, donc quasi-projectif (II, 5.3.4, (i)), et il suffit d'appliquer
(\hyperref[III.4.4.3-fr]{4.4.3}).
\end{proof}

